\documentclass[UTF8, 12pt, fontset=fandol, oneside]{ctexart}
\usepackage{researchnotes}

\graphicspath{{tex/强化学习/figures/math-17/}}

\title{强化学习：应用与案例研究}
\author{}
\date{}

\begin{document}

\maketitle

\section*{第 16 章\quad 应用与案例研究}

在本章中，我们给出强化学习的几个案例研究。
其中若干案例是具有潜在经济意义的重要应用。
其中一个，Samuel 的跳棋程序，主要具有历史兴趣。
我们的介绍旨在说明真实应用中出现的一些权衡和问题。
例如，我们强调领域知识怎样被纳入问题的表述和求解之中。
我们还突出表示问题，因为这些问题常常是应用能否成功的关键。
这些案例研究中使用的算法，有些比本书其余部分介绍过的算法复杂得多。
强化学习的应用仍远未成为常规工作，而且通常既需要技艺，也需要科学。
使应用变得更容易、更直接，是当前强化学习研究的目标之一。

\subsection*{16.1\quad TD-Gammon}

迄今为止，强化学习最令人印象深刻的应用之一，是 Gerald Tesauro 在双陆棋游戏上所做的工作（Tesauro，1992，1994，1995，2002）。
Tesauro 的程序 TD-Gammon 只需要很少的双陆棋知识，却学会了极高水平的对弈，接近世界最强大师的水平。
TD-Gammon 中的学习算法，是 TD($\lambda$) 算法与非线性函数近似的一种直接组合；后者使用多层人工神经网络（ANN），并通过反向传播 TD 误差进行训练。

双陆棋是一个重要游戏，意思是它在全世界都有玩家，有许多锦标赛，也有定期举行的世界冠军赛。
它部分上是一个机会游戏，也是用来下注大量金钱的流行载体。
职业双陆棋棋手的数量，很可能比职业国际象棋棋手还多。
游戏在一个有 24 个位置的棋盘上进行，这些位置称为点位，双方各有 15 枚白子和 15 枚黑子。
下一页右侧给出了游戏早期的一个典型局面，这是从白方视角看到的。
这里白方刚掷出骰子，得到 5 和 2。
这意味着他可以把一枚棋子移动 5 步，并把一枚棋子（可能是同一枚）移动 2 步。
例如，他可以把两枚棋子从 12 点移动出去，一枚移到 17 点，另一枚移到 14 点。
白方的目标是把自己的所有棋子推进到棋盘最后一象限，也就是 19--24 点，然后移出棋盘。
先移出所有棋子的玩家获胜。
一个复杂之处在于，棋子在彼此相向经过时会相互作用。
例如，如果轮到黑方行动，他可以利用骰子点数 2，把一枚棋子从 24 点移动到 22 点，从而“击中”那里的白子。
被击中的棋子会被放到棋盘中央的“栏”上（这里我们已经看到一枚先前被击中的黑子），并从那里重新进入竞赛，从头开始。
然而，如果一个点位上有两枚棋子，那么对手不能移动到这个点位；这些棋子受到保护，不会被击中。
因此，白方不能用他的 5--2 骰子移动 1 点上的任一棋子，因为它们可能到达的点位都被成组的黑子占据。
形成连续的占据点位来阻挡对手，是这个游戏的一种基本策略。

\begin{figure}[htbp]
\centering
\includegraphics[width=0.55\textwidth]{backgammon-position.png}
\caption*{\textbf{图：}一个双陆棋局面。}
\end{figure}

双陆棋还包含若干进一步的复杂性，但以上描述给出了基本思想。
棋盘上有 30 枚棋子和 24 个可能位置；如果把栏和已经移出棋盘的位置也算入，则有 26 个位置。
由此应当清楚，可能的双陆棋局面数量极其巨大，远远超过任何物理上可实现计算机中可拥有的存储单元数。
从每个局面可走的着法数量也很大。
对一次典型掷骰，可能有 20 种不同下法。
在考虑未来着法时，例如对手的回应，还必须考虑可能出现的骰子点数。
结果是，博弈树具有约 400 的有效分支因子。
这太大了，以至于无法有效使用那些在国际象棋和跳棋等游戏中已证明很有效的传统启发式搜索方法。

另一方面，这个游戏很适合 TD 学习方法的能力。
虽然这个游戏具有高度随机性，但游戏状态的完整描述在所有时刻都是可得的。
游戏沿一系列着法和局面展开，直到最终以一方获胜而终止。
结果可以解释为需要预测的最终奖励。
另一方面，我们到目前为止描述的理论结果不能直接有用地应用到这个任务。
状态数太大，无法使用查找表，而对手又是一个不确定性和时间变异性的来源。

TD-Gammon 使用了非线性形式的 TD($\lambda$)。
任一状态（棋盘局面）$s$ 的估计价值 $\hat v(s,\mathbf{w})$，意在估计从状态 $s$ 开始获胜的概率。
为此，所有时间步的奖励都定义为零，只有游戏获胜的时间步除外。
为了实现价值函数，TD-Gammon 使用了一个标准多层 ANN，很像下一页右侧所示的那个。
真实网络在最终层还有两个额外单元，用一种称为“gammon”或“backgammon”的特殊方式估计每个玩家获胜的概率。
该网络由一层输入单元、一层隐藏单元和一个最终输出单元组成。
网络的输入是双陆棋局面的表示，输出则是该局面价值的估计。

\begin{figure}[htbp]
\centering
\includegraphics[width=0.58\textwidth]{figure-16-1.png}
\caption*{\textbf{图 16.1：}TD-Gammon 的人工神经网络。}
\end{figure}

在 TD-Gammon 的第一个版本 TD-Gammon 0.0 中，双陆棋局面以一种相对直接、只涉及少量双陆棋知识的方式表示给网络。
不过，它确实涉及关于 ANN 怎样工作、以及信息怎样最好地呈现给 ANN 的大量知识。
注意 Tesauro 选择的精确表示是很有启发性的。
网络共有 198 个输入单元。
对双陆棋棋盘上的每个点位，四个单元表示该点位上的白子数量。
如果没有白子，那么全部四个单元的值都是零。
如果有一枚棋子，那么第一个单元取值为 1。
这编码了“暴露子”的基本概念，也就是可被对手击中的棋子。
如果有两枚或更多棋子，那么第二个单元被设为 1。
这编码了“占据点”的基本概念，也就是对手不能落子的点位。
如果恰好有三枚棋子，那么第三个单元被设为 1。
这编码了“单个备用子”的基本概念，也就是除形成占据点的两枚棋子之外额外的一枚棋子。
最后，如果有多于三枚的棋子，那么第四个单元被设为与超过三枚的额外棋子数成比例的值。
令 $n$ 表示该点位上的棋子总数；若 $n>3$，则第四个单元取值为 $(n-3)/2$。
这编码了给定点位上“多个备用子”的线性表示。

在 24 个点位上，每个点位对白方有四个单元，对黑方也有四个单元，这就共有 192 个单元。
另外两个单元编码栏上白子和黑子的数量（每个取值为 $n/2$，其中 $n$ 是栏上棋子数），还有两个单元编码黑子和白子已经成功移出棋盘的数量（这些取值为 $n/15$，其中 $n$ 是已经移出的棋子数）。
最后，两个单元以二值方式指示该白方走还是黑方走。
这些选择背后的一般逻辑应当很清楚。
基本上，Tesauro 试图以直接的方式表示局面，同时让单元数量保持相对较小。
他为每个概念上不同且似乎可能相关的量提供一个单元，并把它们缩放到大致相同的范围；在这里，这个范围是 0 到 1 之间。

给定一个双陆棋局面的表示，网络按标准方式计算其估计价值。
从一个输入单元到一个隐藏单元的每条连接，都对应一个实值权重。
来自每个输入单元的信号乘以相应权重，并在隐藏单元处求和。
隐藏单元 $j$ 的输出 $h(j)$ 是加权和的非线性 sigmoid 函数：
\[
h(j)=\sigma\left(\sum_i w_{ij}x_i\right)
=\frac{1}{1+e^{-\sum_i w_{ij}x_i}},
\]
其中 $x_i$ 是第 $i$ 个输入单元的值，$w_{ij}$ 是它到第 $j$ 个隐藏单元的连接权重（网络中的所有权重合在一起构成参数向量 $\mathbf{w}$）。
sigmoid 的输出总在 0 和 1 之间，并且可以自然地解释为基于证据累加得到的概率。
从隐藏单元到输出单元的计算完全类似。
从一个隐藏单元到输出单元的每条连接都有单独的权重。
输出单元形成加权和，然后让它通过同一个 sigmoid 非线性函数。

TD-Gammon 使用了第 12.2 节描述的 TD($\lambda$) 算法的半梯度形式，其中梯度由误差反向传播算法计算（Rumelhart、Hinton 和 Williams，1986）。
回忆这一情形的一般更新规则为
\begin{equation}
\mathbf{w}_{t+1}\doteq \mathbf{w}_t+\alpha\left[R_{t+1}+\gamma\hat v(S_{t+1},\mathbf{w}_t)-\hat v(S_t,\mathbf{w}_t)\right]\mathbf{z}_t,
\tag{16.1}
\end{equation}
其中 $\mathbf{w}_t$ 是所有可修改参数的向量（在此情形中，是网络权重），$\mathbf{z}_t$ 是资格迹向量，即 $\mathbf{w}_t$ 每个分量各有一个资格迹，并按
\[
\mathbf{z}_t\doteq \gamma\lambda\mathbf{z}_{t-1}+\nabla\hat v(S_t,\mathbf{w}_t)
\]
更新，且 $\mathbf{z}_0\doteq \mathbf{0}$。
该式中的梯度可以通过反向传播过程高效计算。
对于双陆棋应用，$\gamma=1$，而奖励除了获胜时之外始终为零，因此学习规则中的 TD 误差部分通常只是 $\hat v(S_{t+1},\mathbf{w})-\hat v(S_t,\mathbf{w})$，如图 16.1 所示。

为了应用学习规则，我们需要双陆棋对局的来源。
Tesauro 让他的学习型双陆棋棋手同自己对弈，从而获得无穷无尽的对局序列。
为了选择着法，TD-Gammon 考虑它在一次掷骰下大约 20 种可行下法，以及这些下法会导致的相应局面。
所得局面是第 6.8 节讨论过的后状态。
网络被查询，用来估计每个后状态的价值。
然后选择会导致最高估计价值局面的着法。
以这种方式继续下去，由 TD-Gammon 替双方走棋，就可以很容易地产生大量双陆棋对局。
每盘棋都被看作一个幕，状态序列为 $S_0,S_1,S_2,\ldots$。
Tesauro 完全增量地应用非线性 TD 规则（16.1），也就是在每次单独着法之后应用。

网络权重最初被设为较小的随机值。
因此，初始评估完全是任意的。
由于着法是基于这些评估选择的，初始着法不可避免地很差，初始对局常常持续数百步甚至数千步，直到某一方几乎偶然获胜。
然而，在几十盘棋之后，性能迅速改善。

在与自己对弈约 300,000 盘之后，如上所述的 TD-Gammon 0.0 学会了下出大致相当于此前最佳双陆棋计算机程序的水平。
这是一个惊人的结果，因为此前所有高性能计算机程序都使用了大量双陆棋知识。
例如，当时的冠军程序可以说是 Neurogammon；它也是 Tesauro 编写的另一个程序，使用 ANN，但不使用 TD 学习。
Neurogammon 的网络是在一个大型训练语料上训练的，该语料由双陆棋专家给出的示范着法组成；此外，它还从一组专门为双陆棋设计的特征开始。
Neurogammon 是一个高度调校、非常有效的双陆棋程序，曾在 1989 年世界双陆棋奥林匹克赛中决定性获胜。
相反，TD-Gammon 0.0 基本上是在零双陆棋知识下构造出来的。
它能达到 Neurogammon 以及所有其他方法的水平，有力证明了自我对弈学习方法的潜力。

TD-Gammon 0.0 在没有专家双陆棋知识时取得的比赛成功，提示了一个显然的修改：加入专门的双陆棋特征，但保留自我对弈 TD 学习方法。
这产生了 TD-Gammon 1.0。
TD-Gammon 1.0 明显远好于此前所有双陆棋程序，并且只有在人类专家中才遇到严肃竞争。
该程序后来的版本 TD-Gammon 2.0（40 个隐藏单元）和 TD-Gammon 2.1（80 个隐藏单元）又加入了选择性的二层搜索过程。
为了选择着法，这些程序不仅前瞻到会立即产生的局面，也前瞻到对手可能的掷骰和着法。
假定对手总是采取眼下看起来对他最好的着法，就可以计算每个候选着法的期望价值，并选择最优者。
为了节省计算时间，第二层搜索只对第一层之后排名较高的候选着法进行，平均约四到五个。
二层搜索只影响所选择的着法；学习过程仍完全照旧进行。
程序的最终版本 TD-Gammon 3.0 和 3.1 使用 160 个隐藏单元和选择性的三层搜索。
TD-Gammon 说明了学习到的价值函数与决策时搜索的结合，这类似于启发式搜索和 MCTS 方法。
在后续工作中，Tesauro 和 Galperin（1997）探索了轨迹采样方法，作为全宽度搜索的替代；它在把每步思考时间保持在约 5--10 秒这一合理范围内的同时，使实战错误率按很大的数值因子降低（4--6 倍）。

在 1990 年代，Tesauro 能够让他的程序同世界级人类棋手进行相当数量的对局。
结果概要见表 16.1。

\begin{table}[htbp]
\centering
\caption*{\textbf{表 16.1：}TD-Gammon 结果概要。}
\begin{tabular}{l c c l l}
\toprule
程序 & 隐藏单元 & 训练对局 & 对手 & 结果 \\
\midrule
TD-Gammon 0.0 & 40 & 300,000 & 其他程序 & 并列最佳 \\
TD-Gammon 1.0 & 80 & 300,000 & Robertie、Magriel 等 & $-13$ 分 / 51 盘 \\
TD-Gammon 2.0 & 40 & 800,000 & 多位大师 & $-7$ 分 / 38 盘 \\
TD-Gammon 2.1 & 80 & 1,500,000 & Robertie & $-1$ 分 / 40 盘 \\
TD-Gammon 3.0 & 80 & 1,500,000 & Kazaros & $+6$ 分 / 20 盘 \\
\bottomrule
\end{tabular}
\end{table}

基于这些结果以及双陆棋大师的分析（Robertie，1992；参见 Tesauro，1995），TD-Gammon 3.0 的棋力看来接近、甚至可能超过世界上最优秀的人类棋手。
Tesauro 在后续文章（Tesauro，2002）中报告了对 TD-Gammon 相对于顶尖人类棋手在走子决策和加倍决策上的广泛 rollout 分析结果。
结论是，在走子决策上，TD-Gammon 3.1 相对于顶尖人类有“压倒性优势”；在加倍决策上则有“轻微优势”。

TD-Gammon 对最优秀人类棋手下棋的方式产生了显著影响。
例如，它学会了用不同于最佳人类棋手惯例的方式来处理某些开局局面。
基于 TD-Gammon 的成功和进一步分析，如今最优秀的人类棋手也按 TD-Gammon 的方式处理这些局面（Tesauro，1995）。
当若干其他受 TD-Gammon 启发的自学 ANN 双陆棋程序广泛可得时，例如 Jellyfish、Snowie 和 GNUBackgammon，它对人类对局的影响被大大加速。
这些程序使 ANN 生成的新知识得以广泛传播，从而使人类锦标赛对局的整体水准大幅提高（Tesauro，2002）。

\subsection*{16.2\quad Samuel 的跳棋程序}

Tesauro 的 TD-Gammon 的一个重要先驱，是 Arthur Samuel（1959，1967）在构造学习下跳棋的程序方面所做的开创性工作。
Samuel 是最早有效使用启发式搜索方法以及我们现在称为时序差分学习的方法的人之一。
他的跳棋程序不仅具有历史兴趣，也是富有启发性的案例研究。
我们强调 Samuel 方法与现代强化学习方法之间的关系，并试图传达 Samuel 使用这些方法的一些动机。

Samuel 在 1952 年最先为 IBM 701 编写了一个下跳棋的程序。
他的第一个学习程序于 1955 年完成，并在 1956 年进行了电视演示。
该程序的后续版本达到良好但非专家级的棋力。
Samuel 被游戏对弈吸引，把它作为研究机器学习的领域，因为游戏比“取自生活”的问题不那么复杂，同时仍允许富有成效地研究启发式过程与学习怎样结合使用。
他选择研究跳棋而不是国际象棋，是因为跳棋相对简单，使人们可以更集中地关注学习。

Samuel 的程序通过从每个当前局面进行前瞻搜索来下棋。
它们使用我们现在称为启发式搜索的方法来决定怎样展开搜索树，以及何时停止搜索。
每次搜索的终端棋盘局面由一个价值函数，或“评分多项式”，用线性函数近似来评估，或“打分”。
在这一点以及其他方面，Samuel 的工作似乎受到 Shannon（1950）建议的启发。
特别地，Samuel 的程序基于 Shannon 的 minimax 过程，从当前局面找出最佳着法。
从已评分的终端局面沿搜索树向后工作，每个局面都被赋予一个由最佳着法产生的局面的分数；这里假定机器总是试图最大化分数，而对手总是试图最小化分数。
Samuel 把这个量称为该局面的“备份分数”。
当 minimax 过程到达搜索树的根，也就是当前局面时，它在这样的假设下给出最佳着法：对手会使用同一个评价准则，只是换到对手的视角。
Samuel 程序的某些版本使用了复杂的搜索控制方法，类似于后来称为“alpha--beta”剪枝的技术（例如参见 Pearl，1984）。

Samuel 使用了两种主要学习方法，其中最简单的一种他称为死记学习。
它只是保存对弈过程中遇到的每个棋盘局面的描述，以及由 minimax 过程确定的该局面的备份价值。
其结果是，如果一个已经遇到过的局面再次作为搜索树的终端局面出现，那么搜索深度会被有效放大，因为该局面的存储价值缓存了此前一次或多次搜索的结果。
一个初始问题是，程序没有被鼓励沿着通往胜利的最直接路径前进。
Samuel 通过在 minimax 分析中每向上备份一层（称为一层 ply）就把局面价值稍微降低一点，给程序一种“方向感”。
他说：“如果程序现在面对一组棋盘局面，其分数只因 ply 数而不同，那么它会自动作出最有利的选择：获胜时选择低 ply 的替代方案，落败时选择高 ply 的替代方案”（Samuel，1959，第 80 页）。
Samuel 发现这种类似折扣的技术对成功学习是必不可少的。
死记学习带来缓慢但持续的改进，且对开局和残局最有效。
他的程序通过同自己下许多盘棋、同各种人类对手下棋，以及在监督学习模式下从棋谱对局学习，成为一个“高于平均水平的新手”。

死记学习以及 Samuel 工作的其他方面，都强烈暗示了时序差分学习的基本思想：一个状态的价值应当等于很可能随后出现的状态的价值。
Samuel 在他的第二种学习方法中最接近这一思想；该方法是用于修改价值函数参数的“通过泛化学习”过程。
Samuel 的方法在概念上与很久以后 Tesauro 在 TD-Gammon 中使用的方法相同。
他让自己的程序与另一个版本的自身进行许多盘对局，并在每次着法之后执行一次更新。
Samuel 更新的思想由图 16.2 中的备份图暗示。
每个空心圆表示程序下一步行动的局面，即轮到己方走的局面；每个实心圆表示对手下一步行动的局面。
每一方各走一步之后，会得到第二个轮到己方走的局面；随后就对每个轮到己方走的局面的价值进行更新。
更新目标是从第二个轮到己方走的局面发起的搜索所得的 minimax 价值。
因此，总体效果是在一个真实事件的完整着法上进行备份，然后再在可能事件上进行一次搜索，如图 16.2 所示。
出于计算原因，Samuel 的实际算法比这复杂得多，但这是基本思想。

Samuel 没有包含显式奖励。
相反，他固定了最重要特征的权重，即棋子优势特征；该特征衡量程序相对于对手拥有多少棋子，对王棋赋予更高权重，并加入一些细化项，使得在获胜时交换棋子比落败时交换棋子更好。
因此，Samuel 程序的目标是改进自己的棋子优势；在跳棋中，这与获胜高度相关。

\begin{figure}[htbp]
\centering
\includegraphics[width=0.62\textwidth]{figure-16-2.png}
\caption*{\textbf{图 16.2：}Samuel 跳棋程序的备份图。}
\end{figure}

然而，Samuel 的学习方法可能遗漏了一个健全时序差分算法的基本部分。
时序差分学习可以看成一种使价值函数与自身一致的方法，而这一点在 Samuel 的方法中清楚可见。
但还需要一种方式，把价值函数固定到状态的真实价值上。
我们通过奖励以及通过折扣或给终止状态赋固定值来强制实现这一点。
但是 Samuel 的方法不包含奖励，也没有对游戏终端局面进行特殊处理。
正如 Samuel 本人指出的，他的价值函数可能仅仅通过给所有局面一个常数值而变得一致。
他希望通过给棋子优势项一个很大的不可修改权重来阻止这种解。
但是，虽然这可能降低找到无用评价函数的可能性，却并不能禁止它们。
例如，仍然可以通过设置可修改权重以抵消不可修改权重的作用，从而达到一个常数函数。

由于 Samuel 的学习过程没有被约束为寻找有用的评价函数，它本应有可能随着经验增加而变差。
事实上，Samuel 报告说，他在大量自我对弈训练过程中观察到了这种情况。
为了让程序重新改进，Samuel 不得不介入，把绝对值最大的权重重新置为零。
他的解释是，这种剧烈干预使程序脱离了局部最优；但另一种可能是，它使程序脱离了那些虽一致却与赢棋或输棋几乎无关的评价函数。

尽管存在这些潜在问题，使用泛化学习方法的 Samuel 跳棋程序仍达到“高于平均水平”的棋力。
相当好的业余对手把它描述为“狡猾但可以击败”（Samuel，1959）。
与死记学习版本相反，这个版本能够发展出良好的中局能力，但在开局和残局上仍然较弱。
该程序还包含一种能力，即搜索一组特征，以找出那些对形成价值函数最有用的特征。
后来的一个版本（Samuel，1967）在搜索过程中加入了改进，例如 alpha--beta 剪枝，广泛使用一种称为“棋谱学习”的监督学习模式，并使用称为签名表的分层查找表（Griffith，1966）来表示价值函数，而不是使用线性函数近似。
这个版本比 1959 年的程序学得好得多，但仍没有达到大师水平。
Samuel 的跳棋程序被广泛认为是人工智能和机器学习中的一项重大成就。

\subsection*{16.3\quad Watson 的 Daily-Double 下注}

IBM Watson\footnote{IBM Corp. 的注册商标。} 是 IBM 研究人员团队开发的系统，用来参加流行电视问答节目 Jeopardy!\footnote{Jeopardy Productions Inc. 的注册商标。}
它在 2011 年通过在一场表演赛中击败人类冠军而成名。
虽然 Watson 所展示的主要技术成就是它能够在广泛的一般知识领域内快速而准确地回答自然语言问题，但它在 Jeopardy! 中的获胜表现也依赖于针对游戏关键部分的复杂决策策略。
Tesauro、Gondek、Lechner、Fan 和 Prager（2012，2013）改造了上面描述的 Tesauro 的 TD-Gammon 系统，以创建 Watson 在著名的对人类冠军获胜表现中用于“Daily-Double”（DD）下注的策略。
这些作者报告说，这一下注策略的有效性远超人类玩家在现场比赛中所能做到的水平；并且它与其他高级策略一起，是 Watson 令人印象深刻的获胜表现的重要贡献因素。
这里我们只关注 DD 下注，因为它是 Watson 中最依赖强化学习的组成部分。

Jeopardy! 由三名选手进行；他们面对一块棋盘，上面显示 30 个方格，每个方格隐藏一条线索并具有一个美元价值。
这些方格排列成六列，每列对应一个不同类别。
一名选手选择一个方格，主持人读出该方格的线索，每名选手都可以选择按响抢答器来回应线索。
第一个抢答的选手可以尝试回答该线索。
如果这名选手回答正确，他的分数增加该方格的美元价值；如果回答不正确，或没有在五秒内回答，则他的分数减少相同数额，而其他选手有机会抢答同一线索。
一个或两个方格（取决于游戏当前轮次）是特殊的 DD 方格。
选择到这些方格之一的选手获得独占机会来回答该方格线索，并且必须在线索揭示之前决定下注多少。
下注额必须大于五美元，但不能超过该选手当前分数。
如果选手正确回答 DD 线索，他的分数增加下注额；否则，分数减少下注额。
每场比赛结束时有一个“Final Jeopardy”（FJ）轮次，每名选手写下一个秘密下注额，然后在线索读出之后写下答案。
三轮比赛之后分数最高的选手获胜；这里一轮指揭示全部 30 条线索。
游戏还有许多其他细节，但这些已经足以理解 DD 下注的重要性。
胜负常常取决于选手的 DD 下注策略。

每当 Watson 选择到一个 DD 方格时，它通过比较动作价值 $\hat q(s,bet)$ 来选择下注额；该动作价值为每个以整美元计的合法下注额估计从当前游戏状态 $s$ 获胜的概率。
除了下面描述的一些风险降低措施外，Watson 选择动作价值最大的下注额。
动作价值在需要下注决策时计算，所用的是两类在任何现场比赛开始前已经学得的估计。
第一类估计是选择每个合法下注额后会产生的后状态价值（第 6.8 节）。
这些估计来自一个状态价值函数 $\hat v(\cdot,\mathbf{w})$，其参数为 $\mathbf{w}$，并对 Watson 从任意游戏状态获胜的概率给出估计。
第二类用于计算动作价值的估计，给出“类别内 DD 置信度” $p_{DD}$，它估计 Watson 正确回答尚未揭示的 DD 线索的可能性。

Tesauro 等人使用上面描述的 TD-Gammon 强化学习方法来学习 $\hat v(\cdot,\mathbf{w})$：这是一个直接组合，即使用多层 ANN 的非线性 TD($\lambda$)，权重 $\mathbf{w}$ 在许多模拟游戏中通过反向传播 TD 误差来训练。
状态以专为 Jeopardy! 设计的特征向量表示给网络。
特征包括三名玩家的当前分数、还剩多少 DD、剩余线索的总美元价值，以及与游戏剩余进程相关的其他信息。
不同于通过自我对弈学习的 TD-Gammon，Watson 的 $\hat v$ 是在数百万盘模拟游戏中学得的，对手是精心构造的人类玩家模型。
类别内置信度估计以 Watson 在当前类别中先前回答正确和回答错误的次数 $(r,w)$ 为条件。
对 $(r,w)$ 的依赖关系，是根据 Watson 在数千个历史类别上的实际正确率估计出来的。

有了先前学得的价值函数 $\hat v$ 和类别内 DD 置信度 $p_{DD}$，Watson 对每个以整美元计的合法下注额如下计算 $\hat q(s,bet)$：
\begin{equation}
\hat q(s,bet)=p_{DD}\times \hat v(S_W+bet,\ldots)
+(1-p_{DD})\times \hat v(S_W-bet,\ldots),
\tag{16.2}
\end{equation}
其中 $S_W$ 是 Watson 的当前分数，而 $\hat v$ 给出 Watson 对 DD 线索回答正确或回答错误之后的游戏状态估计价值。
用这种方式计算动作价值，对应于练习 3.19 的洞见：动作价值是在给定动作下期望的下一状态价值；不过在这里，它是期望的下一后状态价值，因为整个游戏的完整下一状态还取决于下一个方格的选择。

Tesauro 等人发现，通过最大化动作价值来选择下注额会招致“惊人的风险量”；意思是，如果 Watson 对线索的回答碰巧错误，其获胜机会可能遭到灾难性损害。
为了降低错误答案的下行风险，Tesauro 等人通过从 Watson 正确/错误后状态评估值上的标准差中减去一个小分数，对式（16.2）作了调整。
他们还通过禁止某些下注额来进一步降低风险：这些下注额会使错误答案后状态的价值下降到某一限度以下。
这些措施略微降低了 Watson 的获胜期望，但显著降低了下行风险；这不仅体现在每次 DD 下注的平均风险上，在极端风险情形中更是如此，在那些情形中，风险中性的 Watson 会下注其大部分或全部资本。

为什么没有使用 TD-Gammon 的自我对弈方法来学习关键价值函数 $\hat v$？
在 Jeopardy! 中从自我对弈学习不会很好地工作，因为 Watson 与任何人类选手都太不相同。
自我对弈会导致探索一些状态空间区域，而这些区域并不是同人类对手、特别是同人类冠军对局时的典型区域。
此外，不同于双陆棋，Jeopardy! 是一个不完全信息游戏，因为选手无法获得会影响对手对局的全部信息。
特别地，Jeopardy! 选手并不知道对手在各个类别中回答线索有多少信心。
自我对弈会有点像同一个拿着与你相同手牌的人打扑克。

由于这些复杂性，开发 Watson 的 DD 下注策略时，大量工作用于创建良好的人类对手模型。
这些模型并不处理游戏中的自然语言方面，而是对比赛过程中可能发生的事件建立随机过程模型。
统计量来自一个大型的粉丝创建比赛信息档案；该档案覆盖节目从开播到当时的资料。
档案包含诸如线索顺序、选手正确与错误回答、DD 位置，以及近 300,000 条线索上的 DD 和 FJ 下注等信息。
构造了三个模型：平均选手模型（基于全部数据）、冠军模型（基于 100 名最佳选手比赛的统计量），以及超级冠军模型（基于 10 名最佳选手比赛的统计量）。
除了在学习过程中充当对手，这些模型还用于评估学得的 DD 下注策略所带来的收益。
Watson 在模拟中使用基线启发式 DD 下注策略时胜率为 61\%；当它使用学得的价值和一个默认置信度值时，胜率提高到 64\%；而使用现场的类别内置信度时，胜率为 67\%。
考虑到每场比赛中大约只需要 1.5 到 2 次 DD 下注，Tesauro 等人认为这是一个显著改进。

由于 Watson 只有几秒钟可以下注，也只有几秒钟可以选择方格并决定是否抢答，作出这些决策所需的计算时间是一个关键因素。
$\hat v$ 的 ANN 实现使 DD 下注能够足够快地完成，以满足现场比赛的时间约束。
然而，一旦通过改进模拟软件使比赛能够足够快地模拟，在比赛接近结束时，就可以通过对许多蒙特卡洛试验取平均来估计下注价值；在这些试验中，每个下注的后果通过把比赛模拟到结束来确定。
现场比赛中基于蒙特卡洛试验而不是 ANN 来选择残局 DD 下注，显著提高了 Watson 的表现，因为残局中的价值估计误差可能严重影响它的获胜机会。
如果所有决策都通过蒙特卡洛试验来作出，或许会得到更好的下注决策，但鉴于游戏复杂性和现场比赛的时间约束，这根本不可能。

尽管快速而准确地回答自然语言问题的能力是 Watson 的主要成就，但它所有复杂的决策策略都为它令人印象深刻地击败人类冠军作出了贡献。
按照 Tesauro 等人（2012）的说法：
\[
\text{……显然，我们的策略算法达到了超过人类能力的定量精度和实时性能水平。}
\]
这在 DD 下注和残局抢答的情形中尤其如此；在这些情形中，人类根本无法接近 Watson 所执行的精确权益与置信度估计以及复杂决策计算。

\subsection*{16.4\quad 优化存储器控制}

大多数计算机使用动态随机存取存储器（DRAM）作为主存，因为它成本低且容量大。
DRAM 存储器控制器的工作，是高效使用处理器芯片与片外 DRAM 系统之间的接口，以提供高速程序执行所需的高带宽、低延迟数据传输。
存储器控制器需要处理动态变化的读/写请求模式，同时遵守硬件要求的大量时序和资源约束。
这是一个艰巨的调度问题，尤其是在现代处理器中，多个核心共享同一个 DRAM。

\.{I}pek、Mutlu、Mart\'inez 和 Caruana（2008）（另见 Mart\'inez 和 \.{I}pek，2009）设计了一个强化学习存储器控制器，并证明它能够相对于他们研究时常规控制器可达到的水平，显著提高程序执行速度。
他们的动机来自当时已有最先进控制器的局限性；这些控制器使用的策略不能利用过去调度决策的经验，也没有考虑调度决策的长期后果。
\.{I}pek 等人的项目通过模拟完成，但他们在实现所需硬件的详细层面设计了控制器，包括学习算法，使其可以直接在处理器芯片上实现。

访问 DRAM 涉及若干必须按严格时间约束完成的步骤。
DRAM 系统由多个 DRAM 芯片组成，每个芯片包含多个矩形存储单元阵列，这些单元按行和列排列。
每个单元把一个比特存储为电容器上的电荷。
由于电荷会随时间减少，每个 DRAM 单元都需要每隔几毫秒重新充电，即刷新，以防止存储内容丢失。
需要刷新单元，正是 DRAM 被称为“动态”的原因。

每个单元阵列都有一个行缓冲区，用来保存一行比特，这些比特可以被传入阵列的某一行，或从该行传出。
一个 activate 命令会“打开一行”，意思是把命令所指明地址的那一行内容移动到行缓冲区。
在一行已经打开时，控制器可以向单元阵列发出 read 和 write 命令。
每个 read 命令把行缓冲区中的一个字（连续比特的短序列）传输到外部数据总线，每个 write 命令则把外部数据总线中的一个字传输到行缓冲区。
在可以打开另一行之前，必须发出 precharge 命令，把行缓冲区中的（可能已更新的）数据传回到单元阵列中被寻址的行。
之后，另一个 activate 命令可以打开需要访问的新行。
read 和 write 命令称为列命令，因为它们顺序地把比特传入或传出行缓冲区的列；多个比特可以在不重新打开行的情况下传输。
对当前已打开行的 read 和 write 命令，可以比访问另一行更快地执行；访问另一行会涉及额外的行命令，即 precharge 和 activate。
这有时称为“行局部性”。
存储器控制器维护一个存储器事务队列，其中保存来自共享该存储器系统的处理器的存储器访问请求。
控制器必须在遵守大量时序约束的同时，通过向存储器系统发出命令来处理请求。

控制器用于调度访问请求的策略，会对存储器系统性能产生很大影响，例如请求得到满足的平均延迟，以及系统能够达到的吞吐量。
最简单的调度策略按访问请求到达的顺序处理它们：在开始服务下一个请求之前，先发出当前请求所需的全部命令。
但是，如果系统尚未准备好执行其中某个命令，或者执行一个命令会导致资源利用不足（例如，由服务该命令引起的时序约束所致），那么在完成较旧请求之前先开始服务一个较新请求，就是合理的。
策略可以通过重新排序请求获得效率，例如优先处理读请求而不是写请求，或者优先向已经打开的行发出 read/write 命令。
称为 First-Ready, First-Come-First-Serve（FR-FCFS）的策略，把列命令（read 和 write）优先于行命令（activate 和 precharge）；在并列时，则优先选择最旧的命令。
已经证明，在常见条件下，FR-FCFS 在平均存储器访问延迟方面优于其他调度策略（Rixner，2004）。

图 16.3 是 \.{I}pek 等人的强化学习存储器控制器的高层视图。
他们把 DRAM 访问过程建模为一个 MDP，其状态是事务队列的内容，动作是发给 DRAM 系统的命令：precharge、activate、read、write 和 NoOp。
每当动作为 read 或 write 时，奖励信号为 1；否则为 0。
状态转移被认为是随机的，因为系统的下一个状态不仅依赖于调度器的命令，也依赖于系统行为中调度器无法控制的方面，例如访问 DRAM 系统的处理器核心的工作负载。

\begin{figure}[htbp]
\centering
\includegraphics[width=0.65\textwidth]{figure-16-3.png}
\caption*{\textbf{图 16.3：}强化学习 DRAM 控制器的高层视图。调度器是强化学习智能体；其环境由事务队列的特征表示，其动作是发给 DRAM 系统的命令。}
\end{figure}

这个 MDP 的关键，是每个状态中可用动作上的约束。
回忆第 3 章，可用动作集合可以依赖于状态：$A_t\in\mathcal{A}(S_t)$，其中 $A_t$ 是时间步 $t$ 的动作，$\mathcal{A}(S_t)$ 是状态 $S_t$ 中可用动作的集合。
在这一应用中，通过不允许违反时序或资源约束的动作，假定 DRAM 系统的完整性得到保证。
虽然 \.{I}pek 等人没有明确说明，但他们事实上通过为所有可能状态 $S_t$ 预先定义集合 $\mathcal{A}(S_t)$ 来完成这一点。

这些约束解释了为什么该 MDP 有一个 NoOp 动作，也解释了为什么奖励信号除了发出 read 或 write 命令时之外都是 0。
当 NoOp 是某个状态中唯一合法动作时，就会发出 NoOp。
为了最大化存储器系统利用率，控制器的任务是把系统驱动到可以选择 read 或 write 动作的状态：只有这些动作会导致通过外部数据总线发送数据，因此只有它们对系统吞吐量有贡献。
虽然 precharge 和 activate 不产生即时奖励，智能体仍需要选择这些动作，以便之后能够选择带奖励的 read 和 write 动作。

调度智能体使用 Sarsa（第 6.4 节）来学习动作价值函数。
状态由六个整数值特征表示。
为了近似动作价值函数，算法使用线性函数近似，该函数近似通过带哈希的瓦片编码实现（第 9.5.4 节）。
瓦片编码有 32 个铺片，每个铺片存储 256 个动作价值，动作价值以 16 位定点数保存。
探索采用 $\varepsilon$-贪心，$\varepsilon=0.05$。

状态特征包括：事务队列中读请求的数量，事务队列中写请求的数量，事务队列中等待其行被打开的写请求数量，以及事务队列中由发出请求的处理器中最旧者发出、且等待其行被打开的读请求数量。
其他特征依赖于 DRAM 与高速缓存之间如何交互；这里省略细节。
状态特征的选择基于 \.{I}pek 等人对影响 DRAM 性能因素的理解。
例如，根据事务队列中读请求和写请求各有多少来平衡服务读写的速率，可以帮助避免 DRAM 系统与高速缓存交互停滞。
事实上，作者生成了一份相当长的潜在特征列表，然后使用由逐步特征选择引导的模拟，把它们削减到少数几个。

把调度问题表述为 MDP 的一个有趣方面是：输入到瓦片编码、用于定义动作价值函数的特征，不同于用于指定动作约束集合 $\mathcal{A}(S_t)$ 的特征。
瓦片编码输入来自事务队列的内容，而约束集合依赖于许多其他特征，这些特征与整个系统硬件实现必须满足的时序和资源约束有关。
这样，动作约束保证学习算法的探索不会危及物理系统的完整性，同时学习被有效限制在硬件实现的更大状态空间中的一个“安全”区域。

由于这项工作的一个目标是使学习控制器可以在芯片上实现，从而在计算机运行时在线学习，因此硬件实现细节是重要考虑。
设计包含两条五级流水线，用来在每个处理器时钟周期计算和比较两个动作价值，并更新相应的动作价值。
这包括访问存储在芯片上静态 RAM 中的瓦片编码。
对于 \.{I}pek 等人模拟的配置，也就是一个当时高端工作站典型的 4GHz 四核芯片，每个 DRAM 周期对应 10 个处理器周期。
考虑到填充流水线所需的周期，每个 DRAM 周期最多可以评估 12 个动作。
\.{I}pek 等人发现，对任意状态来说，合法命令的数量很少超过这个数；如果总是没有足够时间考虑所有合法命令，性能损失也可以忽略。
这些以及其他巧妙设计细节，使完整控制器和学习算法能够在多处理器芯片上实现。

\.{I}pek 等人通过模拟评估他们的学习控制器，把它同另外三个控制器比较：1. 上面提到的 FR-FCFS 控制器，它产生最佳平均性能；2. 一个按顺序处理每个请求的常规控制器；3. 一个不可实现的理想控制器，称为 Optimistic 控制器，如果有足够需求，它能够通过忽略所有时序和资源约束而维持 100\% 的 DRAM 吞吐量，但除此之外仍对 DRAM 延迟（作为行缓冲区命中）和带宽建模。
他们模拟了九个存储器密集型并行工作负载，这些工作负载由科学应用和数据挖掘应用组成。
图 16.4 显示了四个控制器在九个应用上的性能；性能定义为执行时间的倒数，并归一化到 FR-FCFS 的性能；图中还给出了它们在这些应用上的几何平均性能。
图中标为 RL 的学习控制器，在九个应用上相对于 FR-FCFS 提高了 7\% 到 33\%，平均提高 19\%。
当然，没有任何可实现控制器能够达到 Optimistic 的性能，因为后者忽略所有时序和资源约束；但学习控制器的性能以令人印象深刻的 27\% 缩小了与 Optimistic 上界之间的差距。

\begin{figure}[htbp]
\centering
\includegraphics[width=0.86\textwidth]{figure-16-4.png}
\caption*{\textbf{图 16.4：}四个控制器在 9 个模拟基准应用上的性能。性能归一化到 FR-FCFS，并定义为执行时间的倒数；最右侧是 9 个基准应用上的几何平均性能。}
\end{figure}

由于在芯片上实现学习算法的理由，是让调度策略能够在线适应变化的工作负载，\.{I}pek 等人分析了在线学习相对于先前学得的固定策略的影响。
他们训练

\end{document}
